% Separate post-CP36 candidate. No canonical or checkpoint gate change.
\subsection{An actual third-level arc and the second wild height-two index}

Throughout this subsection $p\ge5$, $q=p^2$, and $N=q-1$.
The integer $\nu$ in B.4 corresponds to actual torsion level
$p^{\nu+1}$. All fields and valuations in the construction are
specified explicitly.

\begin{lemma}[A two-step Newton arc with an actual rational point]
\label{lem:B-RD100-arc}
\status{Candidate; complete proof awaiting independent audits.}
Let $e$ be a positive multiple of $p(p-1)$ with $e>p^2$.
There is a complete discretely valued field $K$ of mixed
characteristic $(0,p)$, residue field $\overline{\mathbf F}_p$,
absolute ramification $e$, and $\mu_p\subset K$, together
with a height-two BT group $G/K$ having a $K$-rational point
of exact order $p^3$. It acquires a good one-dimensional
height-two formal model over a cyclic tame extension $L/K$
of degree $N$.

More precisely, choose a positive integer $A$ such that
\[
 \frac e{p^3}<A<\frac e{p^2},\qquad N\nmid A,
 \quad
 d=eN,\quad b=\frac{eN}{p}-q(p-1)A,\quad J=b+pqA.
                                                        \tag{B-RD100.1}
\]
Such $A$ always exists. For a chosen primitive $N$th root
$\zeta$, the two covariant special eigenvalues of $G$ are
$a,a^p$, where $a=\zeta^{-A}$.
\end{lemma}
\begin{proof}
First put $x_0=e/p^3$. If $x_0<1$, take $A=1$; the upper
inequality follows from $e>p^2$. If $x_0\ge1$, take
$A=\lfloor x_0\rfloor+1$, increasing it by one only if
$N$ divides this integer. Consecutive integers cannot both
be divisible by $N>1$. Thus $x_0<A\le x_0+2<px_0$,
where the last strict inequality uses $p\ge5$ and $x_0\ge1$.
This proves all the required inequalities and nondivisibility.
The integer $b$ is integral because $p\mid e$. Direct algebra gives
\[
\begin{split}
 b-(q-p)A&=\frac Np(e-p^2A)>0,\\
 (q-p)qA-b&=\frac Np(p^3A-e)>0,\\
 d&=b+(p-1)J=pb+pq(p-1)A,\\
 d-J&=(p-1)b+pq(p-2)A>0,\\
 b&\equiv-(p-1)A\pmod N .
\end{split}                                             \tag{B-RD100.2}
\]
In particular $b>0$, $A<qA<J<d$, and
$b<(q-p)qA<(q-p)J$. These are strict inequalities;
no equality polygon is being silently assigned a branch.

Let $k=\overline{\mathbf F}_p$ and $W=W(k)$. Use the
integral $p$-typical formal group $F_u$ over $W[[u]]$
obtained from the original Hopkins--Gross specialization
$v_1=u$, $v_2=1$, $v_i=0$ for $i>2$.
The exact logarithm recursion and coefficient statements are
\[
\begin{gathered}
 f_u(T)=\sum_{j\ge0}l_j(u)T^{p^j},\quad l_0=1,\quad
 p l_j=l_{j-1}u^{p^{j-1}}+l_{j-2},\\
 [p]_u(T)=pT+(1-p^{p-1})uT^p+\sum_{i>p}c_i(u)T^i,\\
 c_i(u)\in(p,u)\quad(p<i<q),\qquad
 c_q(u)\equiv1\pmod{(p,u)}.
\end{gathered}                                           \tag{B-RD100.3}
\]
Here $l_{-1}=0$; the coefficients in degrees $2,\ldots,p-1$
vanish. These are equations (5.3)--(5.5), Proposition~5.7,
and (12.3)--(12.5) of \cite{HopkinsGross1994Bundles}, with
the coefficient of $T^p$ obtained by comparing the degree-$p$
terms of $f_u([p]_uT)=pf_u(T)$. They also appear with their
full specialization argument in Lemma~\ref{lem:B-RD62-arc}.
We use the entire integral power series, including its
unbounded tail, and do not replace it by a polynomial.

In $W[[X]]$ define
\[
 Y_0=X^A,\qquad
 Y_1=[p]_{X^b}(Y_0),\qquad
 Y_2=[p]_{X^b}(Y_1),\qquad
 P(X)=\Phi_{X^b}(Y_2),\quad \Phi_u(T)=[p]_u(T)/T .
                                                        \tag{B-RD100.4}
\]
Every composition is well-defined in the complete power-series
ring. Modulo $p$, the first inequality in (B-RD100.2)
makes the degree-$q$ term uniquely earliest in $Y_1$.
The second makes the degree-$p$ term uniquely earliest in
$Y_2$. Intermediate coefficients in (B-RD100.3) contain $u$
modulo $p$; at each step they are later than the selected term.
Terms of degree greater than $q$ are later even when their
coefficients are units. We obtain the exact leading congruences
\[
 Y_1\equiv X^{qA}+O(X^{qA+1}),\qquad
 Y_2\equiv X^J+O(X^{J+1})\pmod p .
                                                        \tag{B-RD100.5}
\]
In $P$, the $uY_2^{p-1}$ term has order
$b+(p-1)J=d$. All intermediate terms are later, and the
degree-$q$ contribution has order $NJ>d$ because
$b<(q-p)J$. Every term of still higher degree is later.
Therefore
\[
             P(0)=p,\qquad P(X)\equiv X^d+O(X^{d+1})\pmod p.
                                                        \tag{B-RD100.6}
\]
Weierstrass preparation makes $R=W[[X]]/(P)$ the ring
of a degree-$d$ Eisenstein polynomial. It is a complete DVR,
with uniformizer $x$ the image of $X$ and $v_R(p)=d$.
The group $F_{x^b}$ has height-two Honda reduction.
The strict bound $J<d$ ensures that every coefficient
divisible by $p$ is too late to affect (B-RD100.5) on
substituting $x$. Consequently
\[
 v_R(x^A)=A,\qquad v_R(Y_1(x))=qA,\qquad
 v_R(Y_2(x))=J,\qquad [p]^3(x^A)=Y_2(x)P(x)=0 .
                                                        \tag{B-RD100.7}
\]
Since $Y_2(x)\ne0$, this is an actual point of exact order
$p^3$ in the formal group.

We next descend the marked point, not only its valuation.
Take $\zeta\in W^\times$ of order $N$, and let $a=\zeta^{-A}$.
The logarithm recursion proves coefficientwise
\[
 l_j(a^{p-1}u)=a^{p^j-1}l_j(u),\qquad
 [p]_{a^{p-1}u}(a^{-1}T)=a^{-1}[p]_u(T),\qquad
 \Phi_{a^{p-1}u}(a^{-1}T)=\Phi_u(T).
                                                        \tag{B-RD100.8}
\]
For the $l_{j-2}$ summand, the two exponents differ by
$p^j-p^{j-2}$, which is divisible by $N$. The terms of
the recursion thus have the same weight. The resulting
coordinate multiplication $aT$ is integral and has inverse
$a^{-1}T$; the identities are exact over $W[[u]]$.
The congruence for $b$ in (B-RD100.2) gives
$(\zeta X)^b=a^{p-1}X^b$, while
$(\zeta X)^A=a^{-1}X^A$.
Induction through the actual compositions yields
$Y_i(\zeta X)=a^{-1}Y_i(X)$ for $i=0,1,2$ and
$P(\zeta X)=P(X)$.

Thus $\sigma(x)=\zeta x$ defines an automorphism of $R$
of order exactly $N$. Put $L=\operatorname{Frac}(R)$
and $K=L^{\langle\sigma\rangle}$. The actual extension
$L/K$ is cyclic tame of degree $N$, with unchanged residue
field, and $v_K(p)=d/N=e$.
The integral isomorphism $\sigma^*F_{x^b}\to F_{x^b}$
given by $aT$ is a semilinear cocycle. Its $N$fold product
is the identity, and it fixes the marked point since
$a\,\sigma(x^A)=x^A$.
Effective descent at each generic finite level gives
$G/K$ and its $K$-rational point.
The special action on the Honda group is $T\mapsto aT$;
its covariant eigenvalues are $a,a^p$.
This special action need not be faithful when $\gcd(A,N)>1$;
the field action still has order $N$, which is all that the
construction uses.

Finally write $-p=\pi_K^e u$ with $u\in\mathcal O_K^\times$.
Hensel's lemma supplies a $(p-1)$st root of $u$, and
$p-1\mid e$. Hence $K$ contains a $(p-1)$st root of $-p$
and therefore the unique tame cyclotomic extension
$\operatorname{Frac}(W)(\mu_p)$ of degree $p-1$.
This proves $\mu_p\subset K$ and completes the construction.
\end{proof}

\begin{theorem}[Actual third-level character ambiguity]
\label{thm:B-RD100-ambiguity}
\status{Candidate; complete proof awaiting independent audits.}
For every $p\ge5$ and every positive multiple $e$ of
$p(p-1)$ with $e>p^2$, there are height-two BT groups
$M,H$ over a field $K$ as in the preceding lemma such that
\[
                         M[p^3]\simeq H[p^3],           \tag{B-RD100.9}
\]
but their finite special inertia characters differ.
The group $M$ individually acquires supersingular good
reduction; $H$ has a two-step filtration whose rank-one
graded groups acquire good reduction. Thus both satisfy
the stated B.4 monodromy condition. For every $h\ge2$,
$\nu_{\min}(p,e,h)\ge3$ in this range.
\end{theorem}
\begin{proof}
Take the constructed $G$ and actual rational point $P$,
and let $C=\langle P\rangle\simeq\mathbf Z/p^3$.
Its finite flat closure in the good model over
$\mathcal O_L$ and the actual quotient give $G'=G/C$,
which also acquires good reduction. The determinant
argument in B18 and the complete proof of
Theorem~\ref{thm:B-RD62} gives
$\det T_pG=\mathbf Z_p(1)\otimes\psi$, with $\psi$
finite of order dividing $p-1$, trivial on $G_L$, and
$\psi(\sigma)=a^{p+1}$.
This uses the same rank-one crystalline result
\cite[Corollary~2.2.6]{Kisin2006Crystalline} as B62.
No faithful determinant action is required.

The exact multiplication-by-$p^3$ snake sequence is
\[
 0\longrightarrow C\longrightarrow G[p^3]
 \longrightarrow G'[p^3]\longrightarrow C\longrightarrow0 .
                                                        \tag{B-RD100.10}
\]
The first $C$ is a primitive constant rank-one free
$\mathbf Z/p^3$ submodule. Hence its quotient in $G[p^3]$
has character $\chi_p\psi$. After setting
$M=G'\otimes\psi^{-1}$, the actual third truncation fits into
\[
 0\longrightarrow\mu_{p^3}\longrightarrow M[p^3]
 \longrightarrow(\mathbf Z/p^3)(\psi^{-1})
 \longrightarrow0 .                                   \tag{B-RD100.11}
\]
Lemma~\ref{lem:B-RD62-extension} lifts precisely this
extension to a BT group $H$. To recall the effectivity
argument, after killing $\psi$ over a prime-to-$p$ extension,
Kummer identifies the extension classes with
$K'^\times/(K'^\times)^{p^m}$ with the appropriate finite
twist. The transitions are surjective; averaging lifts over
the prime-to-$p$ Galois group preserves surjectivity on
invariants. Correcting compatible representatives by lifted
zero-cochains gives genuinely compatible continuous cocycles.
Their inverse limit defines a free rank-two Tate lattice and
the desired generic BT group. The two graded groups are
$\mu_{p^\infty}$ and
$(\mathbf Q_p/\mathbf Z_p)(\psi^{-1})$.
The whole extension $H$ is not asserted to acquire good
reduction.

Write $\lambda=a^{-1}=\zeta^A$. Isogeny and the finite
twist give the two special eigenvalue multisets
\[
 \begin{array}{c|c}
 M&\{\lambda,\lambda^p\}\\
 H&\{1,\lambda^{p+1}\}.
 \end{array}
 \qquad
 \operatorname{tr}(M)-\operatorname{tr}(H)
             =-(\lambda-1)(\lambda^p-1)\ne0 .           \tag{B-RD100.12}
\]
The nonzero assertion uses exactly $N\nmid A$ and
$\gcd(p,N)=1$. It is independent of whether the finite
special action is faithful. Adding the same constant
height-$h-2$ group preserves the actual truncation isomorphism
and the character difference. Failure at $p^3$, or manuscript
index two, proves the stated lower bound.
\end{proof}

\begin{corollary}[Exact existential mixed threshold at valuation one]
\label{cor:B-RD100-threshold}
\status{Candidate; complete proof awaiting independent audits.}
Among positive multiples $e$ of $p(p-1)$ with $v_p(e)=1$,
there exists an allowed mixed whole-supersingular/filtered
height-two comparison with isomorphic actual $p^3$-torsion
and different finite characters if and only if $e>p^2$.
\end{corollary}
\begin{proof}
The necessary strict inequality is
Proposition~\ref{prop:B-RD76-mixed} at $n=3$ and $v_p(e)=1$.
The theorem constructs actual pairs whenever it holds.
This is an existential statement across fields of the
specified absolute index, not a construction over every
fixed field.
\end{proof}

\begin{theorem}[The exact second wild height-two index]
\label{thm:B-RD100-exact}
\status{Candidate; complete proof awaiting independent audits.}
For every $p\ge5$, in the full B.4 height-two setting
with perfect residue field and $\mu_p\subset K$,
\[
                \nu_{\min}(p,2p(p-1),2)=3.              \tag{B-RD100.13}
\]
In particular $\nu_{\min}(5,40,2)=3$.
\end{theorem}
\begin{proof}
Here $e=2p(p-1)>p^2$ and $v_p(e)=1$, so the theorem
above gives actual $p^3$-torsion failure and the lower bound
three. It may be realized with the particularly simple
choice
\[
 A=1,\quad b=(p-1)(p^2-2),\quad
 J=b+p^3,\quad d=2p(p-1)(p^2-1).                      \tag{B-RD100.14}
\]
All inequalities in the arc proof remain strict.

For the upper bound, consider any allowed actual
$p^4$-torsion comparison at this absolute index.
Lemma~\ref{lem:B-RD76-determinant} detects the total
dimension and finite determinant since $4\ge v_p(e)+2$.
After the completed maximal unramified base change,
which preserves inertia and the absolute index, dimension
zero is recovered from the Brauer lift of semisimplified
residual inertia; dimension two reduces to it by Cartier
duality. In dimension one, two groups filtered by
height-one potentially good factors have their finite
characters determined residually, since $\chi_p\equiv1$
modulo $p$. A mixed comparison would require
$e>p^{2\cdot4-3-1}=p^4$ by
Proposition~\ref{prop:B-RD76-mixed}; this contradicts
$e=2p(p-1)$. Thus only two whole potentially
supersingular groups with common finite determinant remain.

Apply Proposition~\ref{prop:B-RD64-formal} to this actual
comparison with $a=eN=2p(p-1)N$. Its target coefficient
recurrence, starting at $q_0=1$, is
$q_{j+1}=\min\{a+q_j,pq_j+p-1\}$ and gives
\[
 q_1=2p-1,\qquad q_2=2p^2-1,\qquad
 q_3=2p^3-1>2pN=\frac a{p-1}.                       \tag{B-RD100.15}
\]
The strict surplus is $2p-1$.
The actual quotient/cotangent-weight and coefficientwise
logarithm proof of that proposition therefore recovers the
special characters from $p^4$-torsion. It assumes no
source thickness on an isogenous target. This proves
the uniform upper manuscript index three and hence equality.
\end{proof}

\begin{remark}[Construction and scope]
\label{rem:B-RD100-scope}
\status{Candidate; no assertion about two whole good groups is made.}
The lower bound is furnished by the explicit integral
power-series arc, semilinear descent fixing its marked
point, actual finite quotient, and compatible cocycle lift.
It is not inferred from numerical Newton polygons or
character lists. Only one of the two groups is asserted
to acquire good supersingular reduction as a whole.
The auxiliary existential threshold fixes $v_p(e)=1$;
the construction itself allows larger $v_p(e)$ under
its stated divisibility and strict inequality.
The exact index theorem concerns $e=2p(p-1)$ only.

The original Hopkins--Gross equations (5.3)--(5.5),
Proposition~5.7 and its coefficient congruence, and
(12.3)--(12.5), Proposition~12.10 and its full written
proof were freshly reopened and read. The complete B62
arc, quotient and compatible-cocycle proofs and the exact
B64/B76 upper-bound dependencies were reread.
No new page-image inspection or new reading of Kisin
or Bondarko is claimed. The standalone exact supplement
\path{code/reduction_descent/verify_second_wild_trace_arc.py}
passed at 06:09:28.323449 UTC on 20 November 2026 in
0.15055808699980844 seconds. It checks ten symbolic
identities and 5,296 arcs at the 24 primes from five
through101, including 216 corrected choices avoiding
$N\mid A$ and 3,558 nonfaithful special actions.
It checks 100,608 intermediate-coefficient weights,
26,480 tail and $p$-coefficient cutoffs, 14,592 exact
logarithm monomials, and 44 independently inverted
multiplication coefficients through degree $p^2$
at $p=5,7,11,13$. All 24 explicit second-wild
examples and their actual-target upper recurrences
are included. The complete executed program, output
and stdout are retained. No mathematical checker
failure occurred. These computations supplement
the proof; they do not replace any effectivity step.
\end{remark}
